% preamble/theorems.tex
% ------------------------------------
% 概率论定理环境（绿色系 tcolorbox）
% ------------------------------------

% ---- 定义 ----
\newtcbtheorem[number within=section]{mydefinition}{定义}{
	breakable, enhanced,
	colback=white, colframe=GreenMain,
	arc=4pt, boxrule=0.8pt, leftrule=4pt,
	fonttitle=\bfseries\theoremFont,
	fontupper=\normalsize\CJKfamily{theoremCJK},
	left=16pt, right=12pt, top=10pt, bottom=10pt,
	before skip=12pt, after skip=12pt,
	boxsep=4pt, left skip=16pt,
}{def}

% ---- 定理 ----
\newtcbtheorem[use counter from=mydefinition]{mytheorem}{定理}{
	breakable, enhanced,
	colback=white, colframe=GreenMain,
	arc=4pt, boxrule=0.8pt, leftrule=4pt,
	fonttitle=\bfseries\theoremFont,
	fontupper=\normalsize\CJKfamily{theoremCJK},
	left=16pt, right=12pt, top=10pt, bottom=10pt,
	before skip=12pt, after skip=12pt,
	boxsep=4pt, left skip=16pt,
}{th}

% ---- 命题 ----
\newtcbtheorem[use counter from=mydefinition]{proposition}{命题}{
	breakable, enhanced,
	colback=white, colframe=GreenMain,
	arc=4pt, boxrule=0.8pt, leftrule=4pt,
	fonttitle=\bfseries\theoremFont,
	fontupper=\normalsize\CJKfamily{theoremCJK},
	left=16pt, right=12pt, top=10pt, bottom=10pt,
	before skip=12pt, after skip=12pt,
	boxsep=4pt, left skip=16pt,
}{prop}

% ---- 推论 ----
\newtcbtheorem[use counter from=mydefinition]{mycorollary}{推论}{
	breakable, enhanced,
	colback=white, colframe=GreenLight,
	arc=4pt, boxrule=0.8pt, leftrule=4pt,
	fonttitle=\bfseries\theoremFont,
	fontupper=\normalsize\CJKfamily{theoremCJK},
	left=16pt, right=12pt, top=10pt, bottom=10pt,
	before skip=12pt, after skip=12pt,
	boxsep=4pt, left skip=16pt,
}{cor}

% ---- 公理 ----
\newtcbtheorem[use counter from=mydefinition]{myaxiom}{公理}{
	breakable, enhanced,
	colback=white, colframe=GreenDark,
	arc=4pt, boxrule=0.8pt, leftrule=4pt,
	fonttitle=\bfseries\theoremFont,
	fontupper=\normalsize\CJKfamily{theoremCJK},
	left=16pt, right=12pt, top=10pt, bottom=10pt,
	before skip=12pt, after skip=12pt,
	boxsep=4pt, left skip=16pt,
}{ax}

% ---- 引理 ----
\newtcbtheorem[use counter from=mydefinition]{mylemma}{引理}{
	breakable, enhanced,
	colback=white, colframe=GreenAccent,
	arc=4pt, boxrule=0.8pt, leftrule=4pt,
	fonttitle=\bfseries\theoremFont,
	fontupper=\normalsize\CJKfamily{theoremCJK},
	left=16pt, right=12pt, top=10pt, bottom=10pt,
	before skip=12pt, after skip=12pt,
	boxsep=4pt, left skip=16pt,
}{lem}

% ---- 例题 ----
\newtcbtheorem[use counter from=mydefinition]{myexample}{例题}{
	breakable, enhanced,
	colback=white, colframe=GreenLight,
	arc=4pt, boxrule=0.8pt, leftrule=4pt,
	fonttitle=\bfseries\theoremFont,
	fontupper=\normalsize\CJKfamily{theoremCJK},
	left=16pt, right=12pt, top=10pt, bottom=10pt,
	before skip=12pt, after skip=12pt,
	boxsep=4pt, left skip=16pt,
}{ex}

% ---- 例（无色框） ----
\makeatletter
\NewDocumentEnvironment{example}{o}{%
	\par\addvspace{12pt}%
	\refstepcounter{tcb@cnt@mydefinition}%
	\noindent{\bfseries\theoremFont 例~\thetcb@cnt@mydefinition\IfValueT{#1}{：#1}}\par\smallskip
	\normalfont\CJKfamily{theoremCJK}%
}{%
	\par\addvspace{12pt}%
}
\makeatother

% ---- 注记 ----
\newtcbtheorem[use counter from=mydefinition]{myremark}{注记}{
	breakable, enhanced,
	colback=white, colframe=GreenDark,
	arc=4pt, boxrule=0.8pt, leftrule=4pt,
	fonttitle=\bfseries\theoremFont,
	fontupper=\normalsize\CJKfamily{theoremCJK},
	left=16pt, right=12pt, top=10pt, bottom=10pt,
	before skip=12pt, after skip=12pt,
	boxsep=4pt, left skip=16pt,
}{rm}

% ---- 练习（无编号） ----
\NewDocumentEnvironment{exercise}{}{\par\noindent{\bfseries\theoremFont 练习\quad}}{\par}

% ---- 证明环境（楷体，无边框，不参与编号） ----
\RenewDocumentEnvironment{proof}{}{\par\noindent{\bfseries\proofFont 证明.\quad}}{\hfill$\square$\par\bigskip}

% ---- 解答环境 ----
\newenvironment{solution}{%
	\par\proofFont
	\noindent\textbf{解答:}\quad
}{%
	\hfill $\square$\par
}

% ---- 行距 ----
\setstretch{1.2}
